\documentclass[11pt]{article}
\usepackage[a4paper, top=18mm, bottom=18mm, left=20mm, right=20mm]{geometry}
\usepackage[T2A]{fontenc}
\usepackage[utf8]{inputenc}
\usepackage[ukrainian]{babel}
\usepackage{graphicx}
\usepackage[export]{adjustbox}
\usepackage{amsmath}
\usepackage{amssymb}
\graphicspath{{/app/Quiz/Scripts/2023/ProgAll/15BinTree/}{/app/Quiz/Scripts/2021/Mod2021_3/BinTree/}{/app/Quiz/Scripts/2020/Mod2020_4/BinTree/}}
\setlength{\parindent}{0pt}
\setlength{\parskip}{6pt}
\pagestyle{empty}
\begin{document}
%begin
{\large\textbf{Контрольна робота}}

\textbf{Варіант \No\,1}\hfill \textit{ПІБ:}\,\underline{\hspace{60mm}}
\vspace{4mm}

%beginitem
1. %goodpath:Scripts/2018/Mod2018_1/03-good.txt
Відмітити все, що є коректним оператором С++:
( 1 ) bool f(int a);

( 2 ) int f(int a,b);

( 3 ) cout<<3.4;

( 4 ) int n; bool n;
\vspace{5mm}
%enditem
%beginitem
2. Що буде виведено за виконання фрагменту коду?

%code
9**0.5*-2
%code
\vspace{5mm}
%enditem
%beginitem
3. Що буде виведено за виконання фрагменту коду?

%code
print(8 > 4 <= 8.0 is True)
%code
\vspace{5mm}
%enditem
%beginitem
4. %goodpath:Scripts/2019/Mod2019_1/Lex/01-good.txt
Відмітити все, що є однією правильною лексемою:
( 1 ) False

( 2 ) '\'

( 3 ) py

( 4 ) 8,3
\vspace{5mm}
%enditem
%beginitem
5. Що буде виведено за виконання фрагменту коду?

a,b,c=6,2,7
def h(a):
global c    a=3
    b*=2
    c=4
    return a+b+c

a,b,c=9,4,3
print(h(a),a,b,c)
\vspace{5mm}
%enditem
%beginitem
6. Що буде виведено за виконання фрагменту коду?

%code
cout << (12 % 12 % 12 % 12);
%code
\vspace{5mm}
%enditem
%beginitem
7. 

\begin{center}
	\adjustimage{max size={0.8\linewidth}{0.8\paperheight}}
	{../15BinTree/trees_exp/178.png}
\end{center}
Указати послідовність міток вершин дерева, зображеного на рис., що утвориться в результаті обходу дерева в прямому порядку. Мітки записувати через пробіл.\par Алгоритм починає роботу з кореня (на рис. це верхня вершина).
%answer:178 прямому
\vspace{5mm}
%enditem
%beginitem
8. %goodpath:Scripts/2019/Mod2019_1/Lex/03-good.txt
Відмітити все, що є коректним оператором С++:
( 1 ) 3.14%2

( 2 ) if a<b: a=0 else: b=0

( 3 ) i,j=j,j

( 4 ) x+3=5
\vspace{5mm}
%enditem
%beginitem
9. Що буде виведено за виконання фрагменту коду?
%code
a, b, c = 6, 8, 7
def h(a, b, c):
    print(a, b, c, end=" ")

h(b=0, c=2, 1)
print(a, b, c)
%code
\vspace{5mm}
%enditem
%beginitem
10. Що буде виведено за виконання фрагменту коду?

print('R', end=' ')
g=[10,11,12,13,14]
print(g.pop(1), end=' ')
print(g)

\vspace{5mm}
%enditem










%end
\newpage

\newpage

%begin
{\large\textbf{Контрольна робота}}

\textbf{Варіант \No\,2}\hfill \textit{ПІБ:}\,\underline{\hspace{60mm}}
\vspace{4mm}

%beginitem
1. Що буде виведено за виконання фрагменту коду?

a,b,c=6,2,7
def h(a):
global c    a=3
    b*=2
    c=4
    return a+b+c

a,b,c=9,4,3
print(h(a),a,b,c)
\vspace{5mm}
%enditem
%beginitem
2. Що буде виведено за виконання фрагменту коду?

%code
print('R', end=' ')
g = ('a','b','c','d','e',0,1,2,3)
print(g[2])
%code

\vspace{5mm}
%enditem
%beginitem
3. %goodpath:Scripts/2018/Mod2018_1/03-good.txt
Відмітити все, що є коректним оператором С++:
( 1 ) bool f(int a);

( 2 ) int f(int a,b);

( 3 ) cout<<3.4;

( 4 ) int n; bool n;
\vspace{5mm}
%enditem
%beginitem
4. Що буде виведено за виконання фрагменту коду?

%code
cout << (!8<4 and !8.0> true and 9<=4.0);
%code
\vspace{5mm}
%enditem
%beginitem
5. Що буде виведено за виконання фрагменту коду?
%code

a,b,c=6,8,7
def h(a,b,c):
    print(a,b,c,end="")

h(b=0,c=2,1)
print(a,b,c)

%code
\vspace{5mm}
%enditem
%beginitem
6. Що буде виведено за виконання фрагменту коду?

x,y,z=0,2,1
x,z,y,z=6,z,y,x
print(x,y,z)
\vspace{5mm}
%enditem
%beginitem
7. %goodpath:Scripts/2019/Mod2019_1/Lex/03-good.txt
Відмітити все, що є коректним оператором С++:
( 1 ) 3.14%2

( 2 ) if a<b: a=0 else: b=0

( 3 ) i,j=j,j

( 4 ) x+3=5
\vspace{5mm}
%enditem
%beginitem
8. %goodpath:Scripts/2018/Mod2018_1/01-good.txt
Відмітити все, що є однією правильною лексемою:
( 1 ) Switch

( 2 ) 0,5

( 3 ) 8.3

( 4 ) 8,3
\vspace{5mm}
%enditem
%beginitem
9. Що буде виведено за виконання фрагменту коду?

%code
print(9**0.5*-2)
%code
\vspace{5mm}
%enditem
%beginitem
10. Що буде виведено за виконання фрагменту коду?

%code
#include <iostream>
using namespace std;

int h(int a){
    int u = 58;
    if (a < 2) 
        return 6;
    if (a > 5)
         u = 2;
    else
         return 7;
    return u;
}

int main(){
    cout << h(-2);
    return 0;
}
%code

\vspace{5mm}
%enditem










%end
\newpage

\newpage

%begin
{\large\textbf{Контрольна робота}}

\textbf{Варіант \No\,3}\hfill \textit{ПІБ:}\,\underline{\hspace{60mm}}
\vspace{4mm}

%beginitem
1. Що буде виведено за виконання фрагменту коду?

%code
cout << 12 % 12 % 12 % 12;
%code
\vspace{5mm}
%enditem
%beginitem
2. %goodpath:Scripts/2019/Mod2019_1/Lex/01-good.txt
Відмітити все, що є однією правильною лексемою:
( 1 ) False

( 2 ) '\'

( 3 ) py

( 4 ) 8,3
\vspace{5mm}
%enditem
%beginitem
3. Що буде виведено за виконання фрагменту коду?
try:
    res = 9**0.5*-2
    if isinstance(res, bool):
        print(f'bool;{res}')
    elif isinstance(res, int):
        print(f'int;{res}')
    elif isinstance(res, float):
        print(f'float;{res:.2f}')
    else:
        print('???')
except:
    print('Z')
\vspace{5mm}
%enditem
%beginitem
4. Що буде виведено за виконання фрагменту коду?
try:
    res = 5/5*5*5
    if isinstance(res, bool):
        print(f'bool;{res}')
    elif isinstance(res, int):
        print(f'int;{res}')
    elif isinstance(res, float):
        print(f'float;{res:.2f}')
    else:
        print('???')
except:
    print('Z')
\vspace{5mm}
%enditem
%beginitem
5. 
Вкажіть значення та тип виразу:
8<True
\vspace{5mm}
%enditem
%beginitem
6. %goodpath:Scripts/2019/Mod2019_1/Lex/03-good.txt
Відмітити все, що є коректним оператором С++:
( 1 ) 3.14%2

( 2 ) if a<b: a=0 else: b=0

( 3 ) i,j=j,j

( 4 ) x+3=5
\vspace{5mm}
%enditem
%beginitem
7. Що буде виведено за виконання фрагменту коду?
try:
    for a in range(6, 6, -3):
        if a<3:
            break
        print(a, end=';')
        if a<3:
            break
    else:
        print('end',end=';')
    print(a)
except:
    print('Z')
\vspace{5mm}
%enditem
%beginitem
8. Що буде виведено за виконання фрагменту коду?
%code
#include <iostream>
using namespace std;

int f(int &x, int &y){
    x = 8;
    y+= 2;
    return x;
}

int main(){
    int a = 9, b = 5;
    a = f(a, a);
    cout << a << ":" << b <<':';
    {
        int a = 1, b = 6;
        cout << ((a<5) && ((b+=1) < 5))
             << a << ":" << b << ":";
    }
    {
        inta = 4;
        cout << a << ':';
    }
    cout << a << ":" << b << endl;
    return 0;
}
%code

\vspace{5mm}
%enditem
%beginitem
9. %goodpath:Scripts/2018/Mod2018_1/01-good.txt
Відмітити все, що є однією правильною лексемою:
( 1 ) Switch

( 2 ) 0,5

( 3 ) 8.3

( 4 ) 8,3
\vspace{5mm}
%enditem
%beginitem
10. Що буде виведено за виконання фрагменту коду?

%code
print('R', end=' ')
i=6
while i<12:
    i+=1
print(i)
%code

\vspace{5mm}
%enditem










%end
\newpage

\newpage

%begin
{\large\textbf{Контрольна робота}}

\textbf{Варіант \No\,4}\hfill \textit{ПІБ:}\,\underline{\hspace{60mm}}
\vspace{4mm}

%beginitem
1. Що буде виведено за виконання фрагменту коду?
a = 6
b = 2
c = 7

def h(a):
global c    a = 3
    b *= 2
    c = 4
    return a + b + c

a = 9
b = 4
c = 3

try:
    print(h(a), a, b, c, sep=';')
except:
    print('Z', a, b, c, sep=';')
\vspace{5mm}
%enditem
%beginitem
2. 
Записати ВИРАЗ, значенням якого є множина, що складається з цілих чисел від 17 до 2020 (включно), що не діляться на 2.\vspace{5mm}
%enditem
%beginitem
3. %goodpath:Scripts/2019/Mod2019_1/Lex/03-good.txt
Відмітити все, що є коректним оператором С++:
( 1 ) 3.14%2

( 2 ) if a<b: a=0 else: b=0

( 3 ) i,j=j,j

( 4 ) x+3=5
\vspace{5mm}
%enditem
%beginitem
4. 
Записати ВИРАЗ, значенням якого є список, що складається з кубів цілих чисел від 17 до 2020 (включно), що не діляться на жодне з чисел 2, 3, записаних в оберненому порядку.\vspace{5mm}
%enditem
%beginitem
5. %goodpath:Scripts/2018/Mod2018_1/01-good.txt
Відмітити все, що є однією правильною лексемою:
( 1 ) Switch

( 2 ) 0,5

( 3 ) 8.3

( 4 ) 8,3
\vspace{5mm}
%enditem
%beginitem
6. 
Написати програму, що запитує у користувача значення дійсних параметрів $x$ та $y$, обчислює для них значення виразу 
$$\frac{\log_{2} x + 21}{(\tg x + 0.8) (y + 60)}$$ 
та повiдомляє введенi данi та результат обчислення.
 Оцінюється правильність та якість коду.

\vspace{5mm}
%enditem
%beginitem
7. %goodpath:Scripts/2019/Mod2019_1/Lex/01-good.txt
Відмітити все, що є однією правильною лексемою:
( 1 ) False

( 2 ) '\'

( 3 ) py

( 4 ) 8,3
\vspace{5mm}
%enditem
%beginitem
8. Що буде виведено за виконання фрагменту коду?

print('R', end=' ')
for a in range(6,6,-3):
    if a<3:
        break
        print(a, end=' ')
    if a<3:
        break
    else:
        print('end', end=' ')
    print(a)
\vspace{5mm}
%enditem
%beginitem
9. Що буде виведено за виконання фрагменту коду?
g=[10,11,12,13,14]; print(g.pop(1),end=' '); print(g)\vspace{5mm}
%enditem
%beginitem
10. 

\begin{center}
	\adjustimage{max size={0.8\linewidth}{0.8\paperheight}}
	{../15BinTree/trees_exp/178.png}
\end{center}
Указати послідовність міток вершин дерева, зображеного на рис., що утвориться в результаті обходу дерева в прямому порядку. Мітки записувати через пробіл.\par Алгоритм починає роботу з кореня (на рис. це верхня вершина).
%answer:178 прямому
\vspace{5mm}
%enditem










%end
\newpage

\newpage

%begin
{\large\textbf{Контрольна робота}}

\textbf{Варіант \No\,5}\hfill \textit{ПІБ:}\,\underline{\hspace{60mm}}
\vspace{4mm}

%beginitem
1. Що буде виведено за виконання фрагменту коду?

%code
print(9**0.5*-2)
%code
\vspace{5mm}
%enditem
%beginitem
2. %goodpath:Scripts/2019/Mod2019_1/Lex/01-good.txt
Відмітити все, що є однією правильною лексемою:
( 1 ) False

( 2 ) '\'

( 3 ) py

( 4 ) 8,3
\vspace{5mm}
%enditem
%beginitem
3. 

\begin{center}
	\adjustimage{max size={0.8\linewidth}{0.8\paperheight}}
	{../BinTree/trees_exp/178.png}
\end{center}
Указати послідовність міток вершин дерева, зображеного на рис., що утвориться в результаті обходу дерева в прямому порядку. Мітки записувати через пробіл.\par Алгоритм починає роботу з кореня (на рис. це верхня вершина).
%answer:178 прямому
\vspace{5mm}
%enditem
%beginitem
4. Що буде виведено за виконання фрагменту коду?

print('R', end=' ')
i=6
while i<12:
    i+=1
print(i)

\vspace{5mm}
%enditem
%beginitem
5. 
Записати ВИРАЗ, значенням якого є множина, що складається з кубів цілих чисел від 17 до 2020 (включно), що не діляться на жодне з чисел 2, 3.\vspace{5mm}
%enditem
%beginitem
6. %goodpath:Scripts/2019/Mod2019_1/Lex/03-good.txt
Відмітити все, що є коректним оператором С++:
( 1 ) 3.14%2

( 2 ) if a<b: a=0 else: b=0

( 3 ) i,j=j,j

( 4 ) x+3=5
\vspace{5mm}
%enditem
%beginitem
7. Що буде виведено за виконання фрагменту коду?
try:
    i = 6
    while i<12:
        i += 1
    print(i, end=';')
except:
    print('ERR')
\vspace{5mm}
%enditem
%beginitem
8. %goodpath:Scripts/2018/Mod2018_1/01-good.txt
Відмітити все, що є однією правильною лексемою:
( 1 ) Switch

( 2 ) 0,5

( 3 ) 8.3

( 4 ) 8,3
\vspace{5mm}
%enditem
%beginitem
9. Що буде виведено за виконання фрагменту коду?

%code
#include <iostream>
using namespace std;

int a = 6, b = 2, c = 7;

int h(){
    inta = 6;
    intb = 9;
    intc = 4;
    return a + b + c;
}

int main(){
    inta = 3;
    intb = 0;
    intc = 8;
    cout << h() << ':';
    cout << a << ':' << b << ':' << c;
    return 0;
}
%code
\vspace{5mm}
%enditem
%beginitem
10. Що буде виведено за виконання фрагменту коду?

%code
a,b,c=6,2,7
def h(a):
global c    a=3
    b*=2
    c=4
    return a+b+c

a,b,c=9,4,3
print(h(a),a,b,c)
%code
\vspace{5mm}
%enditem










%end
\newpage

\newpage

%begin
{\large\textbf{Контрольна робота}}

\textbf{Варіант \No\,6}\hfill \textit{ПІБ:}\,\underline{\hspace{60mm}}
\vspace{4mm}

%beginitem
1. Що буде виведено за виконання фрагменту коду?

%code
a,b,c=6,2,7
def h(a):
global c    a=3
    b*=2
    c=4
    return a+b+c

a,b,c=9,4,3
print(h(a),a,b,c)
%code
\vspace{5mm}
%enditem
%beginitem
2. %goodpath:Scripts/2019/Mod2019_1/Lex/03-good.txt
Відмітити все, що є коректним оператором С++:
( 1 ) 3.14%2

( 2 ) if a<b: a=0 else: b=0

( 3 ) i,j=j,j

( 4 ) x+3=5
\vspace{5mm}
%enditem
%beginitem
3. Що буде виведено за виконання фрагменту коду?

%code
for a in range(10,10+4,-3):
    if a<3:
        continue
    print(a, end=' ')
    a=10
    if a<3:
        break
else:
    print(13, end=' ')
print(a, end=' ')
%code
\vspace{5mm}
%enditem
%beginitem
4. Що буде виведено за виконання фрагменту коду?

%code
print(8>4>8.0>True)
%code
\vspace{5mm}
%enditem
%beginitem
5. Що буде виведено за виконання фрагменту коду?

%code
a = 6
b = 2
c = 7

def h(a):
global c    a=3
    b=2
    c=4
    return a+b+c

a, b, c = 9,4,3
try:
    print(h(a), a, b, c, sep=';')
except:
    print('Z', a, b, c, sep=';')
%code
\vspace{5mm}
%enditem
%beginitem
6. Що буде виведено за виконання фрагменту коду?

%code
print('R', end=' ')
i=12
while i>6:
    i+=-3
print(i)
%code

\vspace{5mm}
%enditem
%beginitem
7. Що буде виведено за виконання фрагменту коду?

%code
def h(a):
    u=58
    if a<2: 
        return 6
    if a>5:
         u=2
    else:
         return 7
    return u

print(h(-2))
%code
\vspace{5mm}
%enditem
%beginitem
8. Що буде виведено за виконання фрагменту коду?

%code
not8<4 and not8.0>True and 9<=4.0
%code
\vspace{5mm}
%enditem
%beginitem
9. %goodpath:Scripts/2018/Mod2018_1/03-good.txt
Відмітити все, що є коректним оператором С++:
( 1 ) bool f(int a);

( 2 ) int f(int a,b);

( 3 ) cout<<3.4;

( 4 ) int n; bool n;
\vspace{5mm}
%enditem
%beginitem
10. %goodpath:Scripts/2019/Mod2019_1/Lex/01-good.txt
Відмітити все, що є однією правильною лексемою:
( 1 ) False

( 2 ) '\'

( 3 ) py

( 4 ) 8,3
\vspace{5mm}
%enditem










%end
\newpage

\newpage

%begin
{\large\textbf{Контрольна робота}}

\textbf{Варіант \No\,7}\hfill \textit{ПІБ:}\,\underline{\hspace{60mm}}
\vspace{4mm}

%beginitem
1. Що буде виведено за виконання фрагменту коду?

def h(a):
    u=58
    if a<2: 
        return 6
    if a>5:
         u=2
    else:
         return 7
    return u

print(h(-2))
\vspace{5mm}
%enditem
%beginitem
2. %goodpath:Scripts/2019/Mod2019_1/Lex/03-good.txt
Відмітити все, що є коректним оператором С++:
( 1 ) 3.14%2

( 2 ) if a<b: a=0 else: b=0

( 3 ) i,j=j,j

( 4 ) x+3=5
\vspace{5mm}
%enditem
%beginitem
3. Що буде виведено за виконання фрагменту коду?

%code
a,b,c=6,2,7
def h(a):
    a=3
    b=2
    c=4
    return a+b+c

a,b,c=9,4,3
print(h(a),a,b,c)
%code
\vspace{5mm}
%enditem
%beginitem
4. 
Написати програму, що запитує у користувача значення дійсних параметрів $x$ та $y$, обчислює для них значення виразу та повiдомляє введенi данi та результат обчислення.
$$\frac{\log_{2} x + 21}{(\tg x + 0.8) (y + 60)}$$ 

\vspace{5mm}
%enditem
%beginitem
5. %goodpath:Scripts/2018/Mod2018_1/01-good.txt
Відмітити все, що є однією правильною лексемою:
( 1 ) Switch

( 2 ) 0,5

( 3 ) 8.3

( 4 ) 8,3
\vspace{5mm}
%enditem
%beginitem
6. Що буде виведено за виконання фрагменту коду?

%code
i=6
while i<12:
    i+=1
print(i)
%code

\vspace{5mm}
%enditem
%beginitem
7. Що буде виведено за виконання фрагменту коду?

%code
def h(a,b):
    c=58
    if a<2:
        return 6
    if a>5:
         c=2
    else: 
        return 7
    return c

print(h(-2,-2))
%code
\vspace{5mm}
%enditem
%beginitem
8. %goodpath:Scripts/2018/Mod2018_1/03-good.txt
Відмітити все, що є коректним оператором С++:
( 1 ) bool f(int a);

( 2 ) int f(int a,b);

( 3 ) cout<<3.4;

( 4 ) int n; bool n;
\vspace{5mm}
%enditem
%beginitem
9. Що буде виведено за виконання фрагменту коду?

def f(n):
    print("f", end="");
    return n

print(f(-2) and f(1))
\vspace{5mm}
%enditem
%beginitem
10. Що буде виведено за виконання фрагменту коду?

%code
cout << (10 % 10 % 10 % 10);
%code
\vspace{5mm}
%enditem










%end
\newpage

\newpage

%begin
{\large\textbf{Контрольна робота}}

\textbf{Варіант \No\,8}\hfill \textit{ПІБ:}\,\underline{\hspace{60mm}}
\vspace{4mm}

%beginitem
1. Що буде виведено за виконання фрагменту коду?
%code
a, b, c = 6, 8, 7
def h(a, b, c):
    print(a, b, c, end=" ")

h(b=0, c=2, 1)
print(a, b, c)
%code
\vspace{5mm}
%enditem
%beginitem
2. Що буде виведено за виконання фрагменту коду?

%code
i=6
while i<12:
    i+=1
print(i)
%code

\vspace{5mm}
%enditem
%beginitem
3. Що буде виведено за виконання фрагменту коду?

def f(n):
    print("f", end="");
    return n

print(f(-2) and f(1))
\vspace{5mm}
%enditem
%beginitem
4. %goodpath:Scripts/2019/Mod2019_1/Lex/01-good.txt
Відмітити все, що є однією правильною лексемою:
( 1 ) False

( 2 ) '\'

( 3 ) py

( 4 ) 8,3
\vspace{5mm}
%enditem
%beginitem
5. Що буде виведено за виконання фрагменту коду?
%code
#include <iostream>
using namespace std;

int f(int &x, int &y){
    x = 8;
    y+= 2;
    return x;
}

int main(){
    int a = 9, b = 5;
    a = f(a, a);
    cout << a << ":" << b <<':';
    {
        int a = 1, b = 6;
        cout << ((a<5) && ((b+=1) < 5)) << ':';
        cout << a << ":" << b << ':';
    }
    {
        inta = 4;
        cout << a << ':';
    }
    cout << a << ":" << b << endl;
    return 0;
}
%code

\vspace{5mm}
%enditem
%beginitem
6. Що буде виведено за виконання фрагменту коду?

%code
print(5/5*5*5)
%code
\vspace{5mm}
%enditem
%beginitem
7. %goodpath:Scripts/2019/Mod2019_1/Lex/03-good.txt
Відмітити все, що є коректним оператором С++:
( 1 ) 3.14%2

( 2 ) if a<b: a=0 else: b=0

( 3 ) i,j=j,j

( 4 ) x+3=5
\vspace{5mm}
%enditem
%beginitem
8. Що буде виведено за виконання фрагменту коду?

%code
print('R', end=' ')
g=[10,11,12,13,14]
del g[6:]
print(g)
%code

\vspace{5mm}
%enditem
%beginitem
9. 
Вкажіть значення та тип виразу:
not8<4 and not8.0>True and 9<=4.0
\vspace{5mm}
%enditem
%beginitem
10. %goodpath:Scripts/2018/Mod2018_1/01-good.txt
Відмітити все, що є однією правильною лексемою:
( 1 ) Switch

( 2 ) 0,5

( 3 ) 8.3

( 4 ) 8,3
\vspace{5mm}
%enditem










%end
\newpage

\newpage

%begin
{\large\textbf{Контрольна робота}}

\textbf{Варіант \No\,9}\hfill \textit{ПІБ:}\,\underline{\hspace{60mm}}
\vspace{4mm}

%beginitem
1. Що буде виведено за виконання фрагменту коду?
%code
if (15 < 15)
    cout << "x";
else
    cout << "c";
%code

\vspace{5mm}
%enditem
%beginitem
2. Що буде виведено за виконання фрагменту коду?
code
_a = 0

class A:
    _a = 1
    
    def __init__(self):
        self._a = 2
        _a = 3
        A._a = 4

obj = A()
try:
    print(_a, end=' ')
    print(obj._a, end=' ')
    print(A._a, end=' ')
    print(obj._a, end=' ')
    print(obj._A__a, end=' ')
except:
    print('ERR')
code
\vspace{5mm}
%enditem
%beginitem
3. %goodpath:Scripts/2019/Mod2019_1/Lex/03-good.txt
Відмітити все, що є коректним оператором С++:
( 1 ) 3.14%2

( 2 ) if a<b: a=0 else: b=0

( 3 ) i,j=j,j

( 4 ) x+3=5
\vspace{5mm}
%enditem
%beginitem
4. %goodpath:Scripts/2018/Mod2018_1/03-good.txt
Відмітити все, що є коректним оператором С++:
( 1 ) bool f(int a);

( 2 ) int f(int a,b);

( 3 ) cout<<3.4;

( 4 ) int n; bool n;
\vspace{5mm}
%enditem
%beginitem
5. Що буде виведено за виконання фрагменту коду?

%code
print('R', end=' ')
g=[10,11,12,13,14]
del g[6:]
print(g)
%code

\vspace{5mm}
%enditem
%beginitem
6. Що буде виведено за виконання фрагменту коду?

%code
def h(a,b):
    c=58
    if a<2:
        return 6
    if a>5:
         c=2
    else: 
        return 7
    return c

print(h(-2,-2))
%code
\vspace{5mm}
%enditem
%beginitem
7. Що буде виведено за виконання фрагменту коду?

%code
i=6
while i<12:
    i+=1
print(i)
%code

\vspace{5mm}
%enditem
%beginitem
8. %goodpath:Scripts/2019/Mod2019_1/Lex/01-good.txt
Відмітити все, що є однією правильною лексемою:
( 1 ) False

( 2 ) '\'

( 3 ) py

( 4 ) 8,3
\vspace{5mm}
%enditem
%beginitem
9. Що буде виведено за виконання фрагменту коду?

%code
print('R', end=' ')
i=12
while i>6:
    i+=-3
print(i)
%code

\vspace{5mm}
%enditem
%beginitem
10. Що буде виведено за виконання фрагменту коду?

%code
print('R', end=' ')
i=6
while i<12:
    i+=1
print(i)
%code

\vspace{5mm}
%enditem










%end
\newpage

\newpage

%begin
{\large\textbf{Контрольна робота}}

\textbf{Варіант \No\,10}\hfill \textit{ПІБ:}\,\underline{\hspace{60mm}}
\vspace{4mm}

%beginitem
1. Що буде виведено за виконання фрагменту коду?
g=[10,11,12,13,14]; del g[6:]; print(g)\vspace{5mm}
%enditem
%beginitem
2. 

\begin{center}
	\adjustimage{max size={0.8\linewidth}{0.8\paperheight}}
	{../BinTree/trees_exp/178.png}
\end{center}
Указати послідовність міток вершин дерева, зображеного на рис., що утвориться в результаті обходу дерева в прямому порядку. Мітки записувати через пробіл.\par Обхід дерева починається з кореня (на рис. це верхня вершина).
%answer:178 прямому
\vspace{5mm}
%enditem
%beginitem
3. %goodpath:Scripts/2018/Mod2018_1/01-good.txt
Відмітити все, що є однією правильною лексемою:
( 1 ) Switch

( 2 ) 0,5

( 3 ) 8.3

( 4 ) 8,3
\vspace{5mm}
%enditem
%beginitem
4. Що буде виведено за виконання фрагменту коду?
try:
    print(6, end=';')
    print(2j<6j, end=';')
    print(7, end=';')
except TypeError:
    print(9, end=';')
except ValueError:
    print(4, end=';')
except:
    print('Z', end=';')
else:
    print(3, end=';')
finally:
    print(0)
\vspace{5mm}
%enditem
%beginitem
5. %goodpath:Scripts/2018/Mod2018_1/03-good.txt
Відмітити все, що є коректним оператором С++:
( 1 ) bool f(int a);

( 2 ) int f(int a,b);

( 3 ) cout<<3.4;

( 4 ) int n; bool n;
\vspace{5mm}
%enditem
%beginitem
6. Що буде виведено за виконання фрагменту коду?

%code
print(5/5*5*5)
%code
\vspace{5mm}
%enditem
%beginitem
7. %goodpath:Scripts/2019/Mod2019_1/Lex/03-good.txt
Відмітити все, що є коректним оператором С++:
( 1 ) 3.14%2

( 2 ) if a<b: a=0 else: b=0

( 3 ) i,j=j,j

( 4 ) x+3=5
\vspace{5mm}
%enditem
%beginitem
8. Що буде виведено за виконання фрагменту коду?

%code
#include <iostream>

int h(int a, int b){
    int c = 58;
    if (a < 2)
        return 6;
    if (a > 5)
         c = 2;
    else 
        return 7;
    return c;
}

int main(){
    std::cout << h(-2, -2);
    return 0;
}
%code
\vspace{5mm}
%enditem
%beginitem
9. Що буде виведено за виконання фрагменту коду?
try:
    res = 8<True
    if isinstance(res, bool):
        print(f'bool;{res}')
    elif isinstance(res, int):
        print(f'int;{res}')
    elif isinstance(res, float):
        print(f'float;{res:.2f}')
    else:
        print('???')
except:
    print('Z')
\vspace{5mm}
%enditem
%beginitem
10. Що буде виведено за виконання фрагменту коду?
def f(n):
    print('f', end='')
    return n

try:
    print(f(-2) and f(1))
except:
    print('ERR')
\vspace{5mm}
%enditem










%end
\newpage
\end{document}


